\section{Finite-scale angular overlap on the original source}
\label{sec:v70-angular-overlap}
A large radial supremum-to-average ratio does not distinguish an
angular birth from an angular loss. The former needs no compensation
when the comparison collar is farther from the critical value. We
separate these two effects without imposing a bound on that ratio.
The separation is a positive weighting of the original source, not a
change of trajectory or a claim that two angular sets are nested.

Fix $0<\chi\le1/10$ and $0<\varepsilon\le\chi/4$. Retain the
physical charts, the exact label $\lambda=(n,k)$, and the common
coarea versions of Section~\ref{sec:v69-comparison-trace}. Write
\[
 L_z(u)=\ell_{z,\lambda}(\sqrt{2u}),\qquad
 b_z(t_z+u)=b_{z,\lambda}^{\varepsilon,\mathrm{clr}}(t_z+u),
 \qquad 0<u<h_\chi.
\]
The label is suppressed only in these two abbreviations. All angular
fractions and all comparisons below refer to this same exact label.
Let $I=(a,b)$ and $J=(c,d)$, where
$0\le a<b\le h_\chi$ and $0\le c<d<h_\chi$. Both intervals
are fixed before the collision limit. Put $|J|=d-c$ and define
\begin{align}
 a_{z,J}(u)&=\frac1{|J|}\int_J
       \min\{1,L_z(v)/L_z(u)\}\,dv\quad\hbox{if }L_z(u)>0,
                                      \label{eq:v70-overlap-weight}\\
 \Delta_{z,J}(u)&=\frac1{|J|}\int_J(L_z(u)-L_z(v))_+\,dv.
                                      \label{eq:v70-angular-loss}
\end{align}
Set $a_{z,J}(u)=1$ when $L_z(u)=0$. This convention changes no
source: $b_z(t_z+u)=0$ at almost every such level. In particular no
clearance guard, positive center value, or collar mass is a denominator.

\begin{lemma}[Maximal scalar overlap]
\label{lem:v70-scalar-overlap}
The weights are measurable and lie in $[0,1]$. Almost everywhere,
\begin{equation}\label{eq:v70-capacity-identity}
 L_z(u)a_{z,J}(u)=\frac1{|J|}\int_J\min\{L_z(u),L_z(v)\}\,dv,
 \qquad L_z(u)(1-a_{z,J}(u))=\Delta_{z,J}(u).
\end{equation}
For each $u,v$ with $L_z(u)>0$, the integrand in
\eqref{eq:v70-overlap-weight} is the largest $a\in[0,1]$ satisfying
$L_z(u)a\le L_z(v)$.
\end{lemma}
\begin{proof}
The last constraint is equivalent to $a\le L_z(v)/L_z(u)$.
Multiplication by $L_z(u)$ and the identity
$x-\min(x,y)=(x-y)_+$ prove both equalities, also at $L_z(u)=0$.
Measurability follows from parameter integration of bounded measurable
functions. These are statements about angular capacities; they do not
assert the existence of a physical orbit joining the two levels.
\end{proof}

Define two positive source densities on the same physical chart by
\begin{equation}\label{eq:v70-local-source-split}
 s_z^{I,J}(t_z+u)=\one_I(u)a_{z,J}(u)b_z(t_z+u),\qquad
 d_z^{I,J}(t_z+u)=\one_I(u)(1-a_{z,J}(u))b_z(t_z+u).
\end{equation}
Thus $s_z^{I,J}+d_z^{I,J}=\one_I b_z$. This is a partition of
weights, not a decomposition into disjoint events. Each weight is a
measurable insertion into the unchanged physical probability space.
Neither source is normalized again.

For positive source families use
$\mathcal H(f)=\limsup_m\sup_{R,\lambda}m^2\|f_{m;\lambda,R}\|_\infty$.
Write
\[
 W_\delta(b)=\min\{1,\delta+C_g\varepsilon^{-1}\sqrt{2b}\},\qquad
 E_\chi(b,d)=\exp\{K_\chi(\sqrt{2b}+\sqrt{2d})\}.
\]
Here $g_z=D_z^\varepsilon(0)$ is allowed to vanish.

\begin{theorem}[Height without a radial-ratio cutoff]
\label{thm:v70-overlap-height}
Let $s^{I,J;\le\delta}$ be the sum of
\eqref{eq:v70-local-source-split} over any subfamily of the selected
physical charts with $g_z\le\delta$, where $0\le\delta\le1$.
Then
\begin{equation}\label{eq:v70-overlap-height}
 \mathcal H(s^{I,J;\le\delta})
 \le C\frac{d}{d-c}(b-a+d)E_\chi(b,d)W_\delta(b).
\end{equation}
No bound on a birth order, an angular persistence certificate, or a
radial moment is required. On each chart the complementary density
satisfies
\begin{equation}\label{eq:v70-loss-density}
 d_z^{I,J}(t_z+u)
 \le\one_I(u)e^{K_\chi\sqrt{2b}}J_z
          \min\{1,g_z+C_g\varepsilon^{-1}\sqrt{2b}\}
          \Delta_{z,J}(u).
\end{equation}
All source identities and inequalities hold on common coarea versions.
\end{theorem}
\begin{proof}
The physical profile and absolute guard comparison give
$b_z(t_z+u)\le e^{K_\chi\sqrt{2b}}J_zW_\delta(b)L_z(u)$
for almost every $u\in I$ on the indicated subfamily. The first
identity in \eqref{eq:v70-capacity-identity} bounds its product
with $a_{z,J}(u)$ by
\[
 e^{K_\chi\sqrt{2b}}J_zW_\delta(b)
            \frac1{d-c}\int_c^d L_z(v)\,dv
 \le E_\chi(b,d)W_\delta(b)\frac{d}{d-c}
                         \widehat q_{z,d;\lambda}.
\]
The last step uses the lower physical-density bound on $(0,d)$;
the collar average is precisely the unguarded physical average of
\eqref{eq:v69-comparison-trace}. At a roof $t$, contributing
centers lie in $[t-b,t-a]$. Sum over those centers and apply
Theorem~\ref{thm:v69-comparison-collar} with interval length $b-a$
and collar width $d$. This proves \eqref{eq:v70-overlap-height}.
The comparison is made to one positive event after summing the
physical source. In particular it incurs neither a word count nor
an exact-label count and makes no regularity assumption on the
selecting insertion. The second identity in
\eqref{eq:v70-capacity-identity} proves \eqref{eq:v70-loss-density}.
The proof uses only bounded measurable angular profiles.
\end{proof}

\begin{corollary}[All-order and later-birth height without angular loss]
\label{cor:v70-no-loss-height}
Fix $H>0$ with $2H<h_\chi$. Consider any subfamily for which
\begin{equation}\label{eq:v70-no-loss-condition}
 L_z(u)\le L_z(v)\quad\hbox{for almost every }
                    (u,v)\in(0,H)\mathbin{\times}(H,2H).
\end{equation}
Its entire original clearance source on $0<F_z-t_z<H$, with
$g_z\le\delta$, has height at most
\begin{equation}\label{eq:v70-no-loss-height}
 6CH\exp\{(2+\sqrt2)K_\chi\sqrt H\}
            \min\{1,\delta+C_g\varepsilon^{-1}\sqrt{2H}\}.
\end{equation}
At $\delta=0$ this is
$6\sqrt2 CC_g\varepsilon^{-1}H^{3/2}
 \exp\{(2+\sqrt2)K_\chi\sqrt H\}$.
Nondecreasing angular occupation on $(0,2H)$ is sufficient for
\eqref{eq:v70-no-loss-condition}; angular sets themselves need not
be nested. No analytic germ at the center is required.
\end{corollary}
\begin{proof}
Under \eqref{eq:v70-no-loss-condition}, Fubini gives
$\Delta_{z,(H,2H)}=0$ almost everywhere on $(0,H)$. Consequently
$s_z^{(0,H),(H,2H)}=b_z$ there. In
\eqref{eq:v70-overlap-height} put $(a,b,c,d)=(0,H,H,2H)$.
The factor $d(b-a+d)/(d-c)$ equals $6H$. The guard estimate
proves the zero-center assertion, including infinitely flat guards.
\end{proof}

The no-loss condition is an exact finite-scale condition on the
physical angular sets, not a claim that every Lorentz word satisfies
it. It can hold even when the radial ratios are unbounded on a
family. For example, the scalar profiles $L(u)=(u/(2H))^p$ on
$(0,2H)$ have ratio $p+1$ on that whole interval and satisfy
\eqref{eq:v70-no-loss-condition} for every $p\ge0$. A profile
which is zero up to a positive radius and then increases is also
included. These examples explain the distinction from a radial-moment
hypothesis; they do not assert a Lorentz realization of every profile.
An angular spike which subsequently disappears contributes to
\eqref{eq:v70-angular-loss}, rather than being silently controlled
by the corollary.
